% =============================================================================
% APPENDIX BP: METHOD-VOTING: Agent-Based Modeling for Institutional Voting
% =============================================================================
% Module: BCM2_BP_VOTING (ABM Voting Methodology)
% Category: METHOD
% Language: English
% Version: 1.0 (February 2026)
% Status: Draft
%
% Dependencies: PPP (METHOD-ABM), BBB (CORE-WHERE), EEE (METHOD-DESIGN)
% Primary Chapter: PSF-2.0 Model (models/PSF-2-0-PAPAL-SUCCESSION/)
%
% Methodological foundation for agent-based simulation of institutional
% voting processes under supermajority rules. Covers 8 computational
% building blocks: QRE voting, Bayesian belief updating, information
% cascades, committed voters, candidate dropout, supermajority dynamics,
% Monte Carlo validation, and pattern-oriented calibration.
%
% =============================================================================

% ============================================================================
% CROSS-REFERENCE MAP
% ============================================================================

\begin{tcolorbox}[colback=blue!5!white, colframe=blue!75!black, title=Cross-Reference Map]
\small
\textbf{Dependencies (this appendix USES):}
\begin{itemize}[nosep]
\item PPP (METHOD-ABM) --- General ABM methodology, agent architecture
\item BBB (CORE-WHERE) --- Parameter sourcing, literature-validated values
\item EEE (METHOD-DESIGN) --- 9-step model design workflow
\item B (CORE-HOW) --- Complementarity $\gamma_{ij}$ interactions
\end{itemize}

\textbf{Dependents (appendices that USE this):}
\begin{itemize}[nosep]
\item PSF-2.0 Conclave ABM --- Primary application (models/PSF-2-0-PAPAL-SUCCESSION/)
\item BB (DOMAIN-PAPAL-APPOINTMENTS) --- Cardinal appointment analysis
\end{itemize}

\textbf{Level 5 FULL Integration Cross-References:}
\begin{itemize}[nosep]
\item \textbf{Theory Catalog:} CAT-30 (Voting Theory, Social Choice \& Elite Selection), MS-VT-001 through MS-VT-005
\item \textbf{Case Registry:} CAS-940 (Conclave 2025: Leo XIV/Prevost)
\item \textbf{Parameter Registry:} PAR-VTM-001 through PAR-VTM-008
\item \textbf{Working Paper:} WP-2026-001 (outputs/papers/WP\_2026\_001\_PSF-2.0\_nber.md)
\item \textbf{Paper YAML:} PAP-prevost2026psf (integration\_level: I5)
\end{itemize}
\end{tcolorbox}

% ============================================================================
% CHAPTER LINKAGE
% ============================================================================

\begin{tcolorbox}[colback=purple!5!white, colframe=purple!75!black, title=Chapter Linkage]
\textbf{Primary Application:} PSF-2.0 Papal Succession Framework
\begin{itemize}[nosep]
\item Conclave ABM implements all 8 building blocks from this appendix
\item Historical validation against 12 papal conclaves (1878--2025)
\end{itemize}

\textbf{Secondary Chapters:}
\begin{itemize}[nosep]
\item Chapter 13 (HIERARCHY) --- Multi-level decision structures
\item Chapter 11 (HOW) --- Complementarity in voting coalitions
\end{itemize}

\textbf{For modelers:} This appendix provides the methodological toolkit for building
voting ABMs in any institutional context (elections, committees, boards).
\end{tcolorbox}

% ============================================================================
% TITLE
% ============================================================================

\chapter{Appendix BP --- METHOD-VOTING: Agent-Based Modeling for Institutional Voting Processes}
\label{app:bp}

% ============================================================================
% ABSTRACT
% ============================================================================

\begin{tcolorbox}[colback=gray!5!white, colframe=gray!75!black, title=Abstract]
This appendix formalizes the methodological building blocks for agent-based
simulation of institutional voting under supermajority rules. We identify eight
computational components---each grounded in established literature---that jointly
enable realistic simulation of sequential multi-round elections: (1) Quantal
Response voting, (2) Bayesian belief updating, (3) information cascades,
(4) committed voters, (5) candidate dropout, (6) supermajority threshold dynamics,
(7) Monte Carlo variance reduction, and (8) pattern-oriented historical validation.
The framework is applied to the PSF-2.0 papal conclave model but generalizes to
any multi-round institutional election with qualified majority requirements.
\end{tcolorbox}

% ============================================================================
% QUICK REFERENCE
% ============================================================================

\begin{tcolorbox}[colback=yellow!5!white, colframe=yellow!75!black, title=Quick Reference]
\small
\renewcommand{\arraystretch}{1.2}
\begin{tabular}{@{}lll@{}}
\toprule
\textbf{Symbol} & \textbf{Name} & \textbf{Definition} \\
\midrule
$\tau$ & Temperature & Rationality parameter in softmax ($\tau \to 0$: deterministic) \\
$b_i^{(r)}$ & Beliefs & Agent $i$'s belief vector over candidates at round $r$ \\
$\eta$ & Learning rate & Weight on new evidence vs.\ prior ($\eta \in [0,1]$) \\
$\theta_c$ & Commitment threshold & Min.\ belief for deterministic voting \\
$\epsilon$ & Exploration rate & Prob.\ of stochastic vote despite commitment \\
$q^*$ & Supermajority & Required vote share for election ($q^* = 2/3$) \\
$\delta$ & Dropout threshold & Min.\ vote share to remain viable candidate \\
$N_{\text{MC}}$ & Monte Carlo runs & Number of stochastic replications per scenario \\
\bottomrule
\end{tabular}
\end{tcolorbox}

% ============================================================================
% SCOPE
% ============================================================================

\begin{tcolorbox}[
    colback=orange!5!white,
    colframe=orange!75!black,
    title={\textbf{Appendix Scope: Objective / In-Scope / Out-of-Scope / Deliverables}},
    fonttitle=\bfseries
]

\textbf{Objective:}
\begin{itemize}[nosep]
    \item Provide formal, literature-grounded methodology for agent-based voting simulations
\end{itemize}

\textbf{In-Scope:}
\begin{itemize}[nosep]
    \item 8 computational building blocks with axioms and literature sources
    \item Formal definitions for each mechanism (equations, parameters, convergence)
    \item Worked example: PSF-2.0 Conclave ABM with calibration results
    \item Generalization principles for other institutional elections
\end{itemize}

\textbf{Out-of-Scope:}
\begin{itemize}[nosep]
    \item Content-specific papal election analysis (see PSF-2.0 model files)
    \item General ABM methodology without voting focus (see Appendix PPP)
    \item Electoral systems without supermajority (simple majority, plurality)
\end{itemize}

\textbf{Deliverables:}
\begin{itemize}[nosep]
    \item Axiom system (VTM-1 to VTM-8)
    \item Parameter table with literature sources
    \item Validation protocol (pattern-oriented, 3 metrics)
\end{itemize}
\end{tcolorbox}


%%%%%%%%%%%%%%%%%%%%%%%%%%%%%%%%%%%%%%%%%%%%%%%%%%%%%%%%%%%%%%%%%%%%%%%%%%%%%%%
\section{The Fundamental Question}
\label{sec:bp-question}
%%%%%%%%%%%%%%%%%%%%%%%%%%%%%%%%%%%%%%%%%%%%%%%%%%%%%%%%%%%%%%%%%%%%%%%%%%%%%%%

\begin{tcolorbox}[colback=red!5!white, colframe=red!75!black, title=Fundamental Question]
\textbf{How can we simulate institutional voting processes where agents form beliefs,
update them based on observed outcomes, and converge to a supermajority---while
reproducing historically observed patterns of duration and winner selection?}
\end{tcolorbox}

Traditional game-theoretic models of voting assume common knowledge, complete
rationality, and simultaneous choice. Real institutional elections---papal
conclaves, UN Security Council votes, corporate board decisions---violate all
three assumptions. Agents have heterogeneous beliefs, bounded rationality, and
vote sequentially with feedback between rounds.

Agent-based modeling provides a natural framework for capturing these features.
Following Epstein's generative paradigm \citep{epstein2006generative}: if we can
grow the observed voting patterns from micro-level behavioral rules, we have
achieved explanatory sufficiency.

The challenge is methodological: \emph{which} behavioral rules, grounded in
\emph{which} literature, with \emph{which} parameters? This appendix answers
that question systematically.


%%%%%%%%%%%%%%%%%%%%%%%%%%%%%%%%%%%%%%%%%%%%%%%%%%%%%%%%%%%%%%%%%%%%%%%%%%%%%%%
\section{The Eight Building Blocks}
\label{sec:bp-blocks}
%%%%%%%%%%%%%%%%%%%%%%%%%%%%%%%%%%%%%%%%%%%%%%%%%%%%%%%%%%%%%%%%%%%%%%%%%%%%%%%

We identify eight computational building blocks, each with a formal axiom,
a literature foundation, and implementation guidance.


% -------------------------------------------------------------------------
\subsection{Building Block 1: Quantal Response Voting (QRE)}
\label{sec:bp-qre}
% -------------------------------------------------------------------------

\begin{tcolorbox}[colback=blue!3!white, colframe=blue!60!black, title=Axiom VTM-1: Quantal Response]
Agents vote probabilistically according to beliefs transformed by a softmax
function with temperature parameter $\tau > 0$:
\begin{equation}
P(\text{vote for } j \mid b_i) = \frac{\exp(b_{ij} / \tau)}{\sum_k \exp(b_{ik} / \tau)}
\label{eq:qre}
\end{equation}
where $b_{ij}$ is agent $i$'s belief about candidate $j$'s suitability.
\end{tcolorbox}

\textbf{Literature foundation:}
\begin{itemize}[nosep]
\item \citet{mcfadden1974conditional}: Conditional logit analysis establishes the
      mathematical form. The softmax function in Eq.~\eqref{eq:qre} is identical
      to McFadden's logit choice probability.
\item \citet{mckelvey1995quantal}: Quantal Response Equilibrium (QRE) provides the
      game-theoretic justification. Players do not best-respond but ``better-respond''---
      higher-utility actions receive more probability mass, parameterized by $\tau$.
\item \citet{train2009discrete}: Comprehensive treatment of discrete choice with
      simulation, covering mixed logit and Monte Carlo integration.
\end{itemize}

\textbf{Parameter calibration:}
\begin{itemize}[nosep]
\item $\tau \to 0$: Deterministic voting (always vote for highest belief) --- too rigid
\item $\tau \to \infty$: Uniform random voting --- no structure
\item $\tau \approx 1.0$: Balanced regime where beliefs matter but surprises occur
\end{itemize}

In the PSF-2.0 conclave ABM, $\tau = 1.0$ was calibrated to reproduce historical
winner accuracy (100\%) while maintaining realistic vote dispersion in early rounds.


% -------------------------------------------------------------------------
\subsection{Building Block 2: Bayesian Belief Updating}
\label{sec:bp-bayes}
% -------------------------------------------------------------------------

\begin{tcolorbox}[colback=blue!3!white, colframe=blue!60!black, title=Axiom VTM-2: Belief Updating]
After each round, agents update beliefs as a weighted average of prior beliefs
and observed vote shares:
\begin{equation}
b_{ij}^{(r+1)} = (1 - \eta) \cdot b_{ij}^{(r)} + \eta \cdot \hat{v}_j^{(r)}
\label{eq:belief-update}
\end{equation}
where $\hat{v}_j^{(r)} = n_j^{(r)} / N$ is the observed vote share of candidate $j$
in round $r$, and $\eta \in [0,1]$ is the learning rate.
\end{tcolorbox}

\textbf{Literature foundation:}
\begin{itemize}[nosep]
\item \citet{degroot1974reaching}: The DeGroot consensus model establishes this exact
      update rule. Agents reach consensus under mild connectivity conditions.
      Our update in Eq.~\eqref{eq:belief-update} is a DeGroot update where the
      ``neighbor'' is the aggregate vote outcome.
\item \citet{chamley2004rational}: Comprehensive treatment of rational social learning
      where agents combine private signals with public information (observed votes).
\item \citet{ambuhl2021belief}: Empirical evidence on how humans actually update
      beliefs about risks---systematic deviations from Bayes' rule that motivate
      our simplified linear update.
\end{itemize}

\textbf{Why linear, not full Bayesian?}
Full Bayesian updating requires specifying a likelihood model for vote generation,
which introduces additional degrees of freedom. The DeGroot linear update is:
(a) empirically well-documented \citep{degroot1974reaching},
(b) computationally efficient, and
(c) sufficient for convergence---under regularity conditions, it produces the same
long-run consensus as full Bayesian updating \citep{chamley2004rational}.


% -------------------------------------------------------------------------
\subsection{Building Block 3: Information Cascades (Bandwagon)}
\label{sec:bp-cascade}
% -------------------------------------------------------------------------

\begin{tcolorbox}[colback=blue!3!white, colframe=blue!60!black, title=Axiom VTM-3: Bandwagon Acceleration]
When a candidate's vote share exceeds the bandwagon threshold $\beta_{\text{thresh}}$,
beliefs accelerate nonlinearly:
\begin{equation}
b_{ij}^{(r+1)} \leftarrow b_{ij}^{(r+1)} \cdot \left(1 + \alpha_1 \cdot \Delta_j + \alpha_2 \cdot \Delta_j^2 \right)
\label{eq:bandwagon}
\end{equation}
where $\Delta_j = \hat{v}_j^{(r)} - \beta_{\text{thresh}}$ is the excess vote share
above the bandwagon threshold, $\alpha_1$ is the linear acceleration, and $\alpha_2$
is the quadratic amplification.
\end{tcolorbox}

\textbf{Literature foundation:}
\begin{itemize}[nosep]
\item \citet{bikhchandani1992theory}: Seminal cascade theory. When public information
      (observed votes) becomes sufficiently strong, agents rationally ignore their
      private signals and follow the crowd. The cascade is \emph{fragile} early but
      \emph{robust} once established---motivating our threshold-based activation.
\item \citet{banerjee1992simple}: Complementary herding model. In sequential decisions,
      rational agents imitate predecessors. Combined with Bikhchandani, this provides
      both the threshold trigger and the acceleration mechanism.
\item \citet{lux1995herd}: Introduces \emph{nonlinear} bandwagon dynamics with
      quadratic amplification in financial markets. Our $\alpha_2 \cdot \Delta^2$
      term directly follows Lux's insight that herding intensifies superlinearly.
\end{itemize}

\textbf{Calibration insight:} The bandwagon threshold $\beta_{\text{thresh}} = 0.33$
corresponds to the critical mass identified in cascade experiments---roughly one-third
of the group must agree before cascades begin \citep{bikhchandani1992theory}.


% -------------------------------------------------------------------------
\subsection{Building Block 4: Committed Voters (Zealots)}
\label{sec:bp-zealot}
% -------------------------------------------------------------------------

\begin{tcolorbox}[colback=blue!3!white, colframe=blue!60!black, title=Axiom VTM-4: Committed Voting]
An agent with belief $b_{ij} > \theta_c$ for some candidate $j$ votes
deterministically for $j$ with probability $(1 - \epsilon)$ and explores
(votes probabilistically via QRE) with probability $\epsilon$:
\begin{equation}
\text{vote}_i = \begin{cases}
j^* = \arg\max_j b_{ij} & \text{with prob. } (1 - \epsilon) \text{ if } b_{ij^*} > \theta_c \\
\text{QRE}(b_i, \tau) & \text{otherwise}
\end{cases}
\label{eq:committed}
\end{equation}
The commitment threshold decreases over rounds: $\theta_c^{(r)} = \max(0.35, \theta_c^{(0)} - 0.02r)$.
\end{tcolorbox}

\textbf{Literature foundation:}
\begin{itemize}[nosep]
\item \citet{mobilia2003single}: A single zealot (committed voter) fundamentally
      alters the long-run dynamics of a voter model---preventing consensus in some
      configurations, accelerating it in others. Our committed voters serve as
      coalition anchors.
\item \citet{galam2007inflexible}: Inflexible minorities can ``break'' democratic
      opinion dynamics. In conclave terms: a small bloc of committed supporters
      prevents early elimination of their candidate, creating a holding pattern
      until institutional pressure (Building Block 6) forces convergence.
\item \citet{masuda2012evolution}: Committed agents act as catalysts for coalition
      formation, disproportionate to their number. This explains the kingmaker
      phenomenon in conclaves.
\end{itemize}

\textbf{Design choice:} The exploration parameter $\epsilon = 0.08$ ensures that even
committed voters occasionally deviate, preventing artificial lock-in. This captures
the empirical observation that even strong supporters occasionally cast ``signal votes''
for alternative candidates.


% -------------------------------------------------------------------------
\subsection{Building Block 5: Candidate Dropout}
\label{sec:bp-dropout}
% -------------------------------------------------------------------------

\begin{tcolorbox}[colback=blue!3!white, colframe=blue!60!black, title=Axiom VTM-5: Candidate Viability]
Candidates whose vote share falls below the dropout threshold $\delta$ suffer
accelerated belief decay:
\begin{equation}
b_{ij}^{(r+1)} \leftarrow b_{ij}^{(r+1)} \cdot \phi_{\text{drop}} \quad \text{if } \hat{v}_j^{(r)} < \delta
\label{eq:dropout}
\end{equation}
where $\phi_{\text{drop}} \in (0,1)$ is the dropout decay factor.
\end{tcolorbox}

\textbf{Literature foundation:}
\begin{itemize}[nosep]
\item This mechanism combines insights from cascade theory \citep{bikhchandani1992theory}
      and the zealot literature \citep{galam2007inflexible}. When support falls below
      a critical mass, the information cascade \emph{reverses}: agents interpret low
      vote share as negative public information and abandon the candidate.
\item In the strategic voting literature (reviewed in \citet{mckelvey1995quantal}),
      this corresponds to the elimination of dominated strategies through iterated
      reasoning---voters stop supporting candidates with no path to victory.
\end{itemize}

\textbf{Calibration:} $\delta = 0.05$ (5\% vote share) and $\phi_{\text{drop}} = 0.5$
ensure that non-viable candidates are eliminated within 2--3 rounds, consistent with
the historical observation that papal conclaves typically narrow to 2--3 serious
candidates after the first few scrutinies.


% -------------------------------------------------------------------------
\subsection{Building Block 6: Supermajority Threshold Dynamics}
\label{sec:bp-supermajority}
% -------------------------------------------------------------------------

\begin{tcolorbox}[colback=blue!3!white, colframe=blue!60!black, title=Axiom VTM-6: Qualified Majority]
The election terminates when any candidate receives vote share $\hat{v}_j \geq q^*$,
where $q^* = 2/3$ is the supermajority threshold. The threshold $q^*$ creates a
strategic tension between the extensive margin (how many candidates) and the
intensive margin (how concentrated support must be).
\end{tcolorbox}

\textbf{Literature foundation:}
\begin{itemize}[nosep]
\item \citet{buchanan1962calculus}: Chapter 15 analyzes qualified majority rules.
      The $2/3$ threshold reflects the constitutional designer's judgment that the
      external costs of a ``wrong'' election (schism, illegitimacy) justify the
      higher decision costs (longer conclave, more rounds).
\item \citet{caplin1988majority}: Proves that the $64\%$-majority rule (approximately
      $2/3$) is welfare-optimal under mild distributional assumptions. The conclave's
      $2/3$ requirement is thus not arbitrary but near-optimal in a Condorcet sense.
\end{itemize}

\textbf{Computational implication:} The $2/3$ threshold means that with 5 candidates,
at least 60\% of the ``belief mass'' must concentrate on one candidate. This is
achievable only through the interaction of Blocks 1--5: QRE concentrates votes,
belief updating creates momentum, cascades amplify it, committed voters anchor
coalitions, and dropout reduces the candidate field.


% -------------------------------------------------------------------------
\subsection{Building Block 7: Monte Carlo Variance Reduction}
\label{sec:bp-montecarlo}
% -------------------------------------------------------------------------

\begin{tcolorbox}[colback=blue!3!white, colframe=blue!60!black, title=Axiom VTM-7: Stochastic Replication]
Each scenario is simulated $N_{\text{MC}}$ times with different random seeds.
The reported outcome is the median over replications:
\begin{equation}
\hat{R} = \text{median}\left\{R^{(1)}, R^{(2)}, \ldots, R^{(N_{\text{MC}})}\right\}
\label{eq:mc}
\end{equation}
where $R^{(k)}$ is the number of rounds to convergence in replication $k$.
\end{tcolorbox}

\textbf{Literature foundation:}
\begin{itemize}[nosep]
\item \citet{robert2004monte}: Standard reference for Monte Carlo statistical methods.
      Variance reduction through replication is a fundamental technique.
\item \citet{grazzini2015estimation}: Simulated minimum distance estimation for
      ergodic ABMs---provides theoretical justification for using multiple stochastic
      runs to estimate ABM parameters when the likelihood function is intractable.
\end{itemize}

\textbf{Design choice:} $N_{\text{MC}} = 5$ balances computational cost with variance
reduction. For the conclave ABM, increasing to $N_{\text{MC}} = 100$ changes median
round estimates by $< 0.5$ rounds, suggesting rapid convergence of the estimator.
Median (not mean) is used for robustness against occasional outlier runs.


% -------------------------------------------------------------------------
\subsection{Building Block 8: Pattern-Oriented Historical Validation}
\label{sec:bp-validation}
% -------------------------------------------------------------------------

\begin{tcolorbox}[colback=blue!3!white, colframe=blue!60!black, title=Axiom VTM-8: Multi-Pattern Validation]
The ABM is validated against $K$ historical cases using three complementary metrics:
\begin{enumerate}[nosep]
\item \textbf{Winner accuracy:} $\text{Acc} = \frac{1}{K}\sum_{k=1}^{K} \mathbb{1}[\hat{w}_k = w_k]$
\item \textbf{Duration error:} $\text{MAE} = \frac{1}{K}\sum_{k=1}^{K} |R_k - \hat{R}_k|$
\item \textbf{Duration correlation:} $r = \text{cor}(R_1, \ldots, R_K; \hat{R}_1, \ldots, \hat{R}_K)$
\end{enumerate}
A model passes validation if $\text{Acc} > 0.80$, $\text{MAE} < 5.0$, and $r > 0.40$.
\end{tcolorbox}

\textbf{Literature foundation:}
\begin{itemize}[nosep]
\item \citet{grimm2005pattern}: Pattern-Oriented Modeling (POM) establishes the
      principle that ABMs should be validated against \emph{multiple} observed patterns
      at \emph{different hierarchical levels}. Our three metrics capture:
      (a) the categorical pattern (who wins),
      (b) the quantitative pattern (how many rounds), and
      (c) the relational pattern (do contested elections take longer?).
\item \citet{windrum2007empirical}: Distinguishes input validation, process validation,
      and output validation. Our approach combines all three: historical papabile scores
      as input, mechanism effects as process, and winner/duration as output.
\end{itemize}

\textbf{State of the art:} \citet{crokidakis2025conclave} validates a 2-parameter
conclave ABM (imitation probability $p$, strategic switching $q$) against historical
data 1939--2025. Our PSF-2.0 ABM achieves higher accuracy with richer behavioral
mechanisms (8 building blocks vs.\ 2 parameters).


%%%%%%%%%%%%%%%%%%%%%%%%%%%%%%%%%%%%%%%%%%%%%%%%%%%%%%%%%%%%%%%%%%%%%%%%%%%%%%%
\section{Axiom System Summary}
\label{sec:bp-axioms}
%%%%%%%%%%%%%%%%%%%%%%%%%%%%%%%%%%%%%%%%%%%%%%%%%%%%%%%%%%%%%%%%%%%%%%%%%%%%%%%

\begin{table}[htbp]
\centering
\caption{Axiom System for Voting ABMs}
\label{tab:bp-axioms}
\renewcommand{\arraystretch}{1.3}
\small
\begin{tabular}{@{}clll@{}}
\toprule
\textbf{Axiom} & \textbf{Name} & \textbf{Key Reference} & \textbf{Equation} \\
\midrule
VTM-1 & Quantal Response & McKelvey \& Palfrey (1995) & Eq.~\eqref{eq:qre} \\
VTM-2 & Belief Updating  & DeGroot (1974) & Eq.~\eqref{eq:belief-update} \\
VTM-3 & Bandwagon & Bikhchandani et al.\ (1992) & Eq.~\eqref{eq:bandwagon} \\
VTM-4 & Committed Voters & Mobilia (2003), Galam (2007) & Eq.~\eqref{eq:committed} \\
VTM-5 & Candidate Dropout & Bikhchandani (1992), reverse cascade & Eq.~\eqref{eq:dropout} \\
VTM-6 & Supermajority & Buchanan \& Tullock (1962) & $q^* = 2/3$ \\
VTM-7 & Monte Carlo & Robert \& Casella (2004) & Eq.~\eqref{eq:mc} \\
VTM-8 & Pattern Validation & Grimm et al.\ (2005) & Acc, MAE, $r$ \\
\bottomrule
\end{tabular}
\end{table}


%%%%%%%%%%%%%%%%%%%%%%%%%%%%%%%%%%%%%%%%%%%%%%%%%%%%%%%%%%%%%%%%%%%%%%%%%%%%%%%
\section{Integration with EBF Framework}
\label{sec:bp-integration}
%%%%%%%%%%%%%%%%%%%%%%%%%%%%%%%%%%%%%%%%%%%%%%%%%%%%%%%%%%%%%%%%%%%%%%%%%%%%%%%

\subsection{Connection to Appendix PPP (METHOD-ABM)}

Appendix PPP provides the \emph{general} ABM methodology for the EBF (agent architecture,
convergence theory, Mean Field limits). This appendix (BP) specializes PPP for
\emph{voting processes}, adding:

\begin{itemize}[nosep]
\item Discrete outcomes (winner selection) vs.\ continuous trajectories (life journeys)
\item Supermajority termination condition vs.\ time-horizon termination
\item Round-based dynamics vs.\ continuous-time update
\item Historical validation via specific elections vs.\ population-level patterns
\end{itemize}

The relationship is:
\[
\text{PPP (general ABM)} \supset \text{BP (voting ABM)} \supset \text{PSF-2.0 (conclave ABM)}
\]

\subsection{Connection to 10C Framework}

The 8 building blocks map to 10C CORE questions:

\begin{table}[htbp]
\centering
\small
\renewcommand{\arraystretch}{1.2}
\begin{tabular}{@{}lll@{}}
\toprule
\textbf{Building Block} & \textbf{10C CORE} & \textbf{Connection} \\
\midrule
VTM-1 (QRE) & WHAT (C) & Utility/belief determines choice \\
VTM-2 (Beliefs) & AWARE (AU) & Awareness of others' votes \\
VTM-3 (Cascades) & HOW (B) & Complementarity in coalition formation \\
VTM-4 (Zealots) & WHO (AAA) & Agent heterogeneity, faction identity \\
VTM-5 (Dropout) & STAGE (AW) & Journey through candidate viability \\
VTM-6 (Supermajority) & HIERARCHY (HI) & Institutional design constraint \\
VTM-7 (Monte Carlo) & WHERE (BBB) & Parameter uncertainty quantification \\
VTM-8 (Validation) & WHEN (V) & Context-specific calibration \\
\bottomrule
\end{tabular}
\end{table}


%%%%%%%%%%%%%%%%%%%%%%%%%%%%%%%%%%%%%%%%%%%%%%%%%%%%%%%%%%%%%%%%%%%%%%%%%%%%%%%
\section{Worked Example: PSF-2.0 Conclave ABM}
\label{sec:bp-example}
%%%%%%%%%%%%%%%%%%%%%%%%%%%%%%%%%%%%%%%%%%%%%%%%%%%%%%%%%%%%%%%%%%%%%%%%%%%%%%%

\subsection{Setup}

The PSF-2.0 Conclave ABM simulates 120 cardinal-agents voting for 5 papabili
candidates over multiple rounds with the following calibrated parameters:

\begin{table}[htbp]
\centering
\caption{PSF-2.0 Calibrated Parameters}
\label{tab:bp-params}
\renewcommand{\arraystretch}{1.2}
\small
\begin{tabular}{@{}llll@{}}
\toprule
\textbf{Parameter} & \textbf{Symbol} & \textbf{Value} & \textbf{Source} \\
\midrule
Temperature & $\tau$ & 1.0 & Calibrated (QRE range 0.5--2.0) \\
Learning rate & $\eta$ & 0.30 & DeGroot (1974), calibrated \\
Bandwagon threshold & $\beta_{\text{thresh}}$ & 0.33 & Bikhchandani et al.\ (1992) \\
Linear acceleration & $\alpha_1$ & 0.4 & Lux (1995), calibrated \\
Quadratic acceleration & $\alpha_2$ & 0.8 & Lux (1995), calibrated \\
Commitment threshold & $\theta_c$ & 0.60 & Galam (2007), calibrated \\
Exploration rate & $\epsilon$ & 0.08 & Masuda (2012), calibrated \\
Dropout threshold & $\delta$ & 0.05 & Empirical (conclave histories) \\
Dropout factor & $\phi_{\text{drop}}$ & 0.50 & Calibrated \\
Monte Carlo runs & $N_{\text{MC}}$ & 5 & Robert \& Casella (2004) \\
\bottomrule
\end{tabular}
\end{table}

\subsection{Results}

Validation against 12 historical papal conclaves (1939--2013):

\begin{table}[htbp]
\centering
\caption{PSF-2.0 Historical Validation Results}
\label{tab:bp-results}
\renewcommand{\arraystretch}{1.2}
\small
\begin{tabular}{@{}lccc@{}}
\toprule
\textbf{Metric} & \textbf{Threshold} & \textbf{Achieved} & \textbf{Status} \\
\midrule
Winner accuracy & $> 0.80$ & $1.00$ (12/12) & $\checkmark$ \\
Mean duration error & $< 5.0$ rounds & $2.08$ rounds & $\checkmark$ \\
Duration correlation & $> 0.40$ & $r = 0.616$ & $\checkmark$ \\
\bottomrule
\end{tabular}
\end{table}

The model reproduces three levels of pattern (following Grimm et al., 2005):
\begin{enumerate}[nosep]
\item \textbf{Categorical:} Correct winner in all 12 conclaves
\item \textbf{Quantitative:} Duration within 2 rounds on average
\item \textbf{Relational:} Contested elections (low papabile dominance) take longer ($r = 0.616$)
\end{enumerate}


%%%%%%%%%%%%%%%%%%%%%%%%%%%%%%%%%%%%%%%%%%%%%%%%%%%%%%%%%%%%%%%%%%%%%%%%%%%%%%%
\section{Generalization: Beyond Papal Conclaves}
\label{sec:bp-generalization}
%%%%%%%%%%%%%%%%%%%%%%%%%%%%%%%%%%%%%%%%%%%%%%%%%%%%%%%%%%%%%%%%%%%%%%%%%%%%%%%

The 8 building blocks generalize to any institutional election with:
\begin{itemize}[nosep]
\item Sequential rounds with public feedback (boards, committees, party leadership)
\item Qualified majority threshold (UN Security Council: 9/15, EU Council: 55\%)
\item Heterogeneous agents with private beliefs (factions, parties, interest groups)
\item Institutional pressure mechanisms (deadlines, recesses, external events)
\end{itemize}

To apply the framework to a new domain:
\begin{enumerate}[nosep]
\item Define agent types and initial belief formation (Building Block 1)
\item Set institutional threshold $q^*$ (Building Block 6)
\item Calibrate learning rate $\eta$ and cascade threshold $\beta$ from historical data (Blocks 2--3)
\item Determine if committed voters/zealots exist in the institution (Block 4)
\item Validate against historical elections using Pattern-Oriented Modeling (Block 8)
\end{enumerate}


%%%%%%%%%%%%%%%%%%%%%%%%%%%%%%%%%%%%%%%%%%%%%%%%%%%%%%%%%%%%%%%%%%%%%%%%%%%%%%%
\section{Summary}
\label{sec:bp-summary}
%%%%%%%%%%%%%%%%%%%%%%%%%%%%%%%%%%%%%%%%%%%%%%%%%%%%%%%%%%%%%%%%%%%%%%%%%%%%%%%

\begin{tcolorbox}[colback=green!5!white, colframe=green!75!black, title=Summary]
\textbf{Appendix BP} provides 8 literature-grounded building blocks for
agent-based simulation of institutional voting:

\begin{enumerate}[nosep]
\item \textbf{QRE:} Probabilistic voting via softmax (McKelvey \& Palfrey, 1995)
\item \textbf{Beliefs:} DeGroot consensus updating (DeGroot, 1974)
\item \textbf{Cascades:} Nonlinear bandwagon (Bikhchandani et al., 1992; Lux, 1995)
\item \textbf{Zealots:} Committed voters as coalition anchors (Mobilia, 2003; Galam, 2007)
\item \textbf{Dropout:} Reverse cascades eliminate non-viable candidates
\item \textbf{Supermajority:} Qualified majority as near-optimal rule (Caplin \& Nalebuff, 1988)
\item \textbf{Monte Carlo:} Stochastic replication for variance reduction
\item \textbf{Validation:} Pattern-oriented multi-metric approach (Grimm et al., 2005)
\end{enumerate}

Applied to the PSF-2.0 Conclave ABM: 100\% winner accuracy, 2.08 mean round error,
$r = 0.616$ duration correlation across 12 historical conclaves (1939--2013).
\end{tcolorbox}


%%%%%%%%%%%%%%%%%%%%%%%%%%%%%%%%%%%%%%%%%%%%%%%%%%%%%%%%%%%%%%%%%%%%%%%%%%%%%%%
\section{Glossary of Symbols}
\label{sec:bp-glossary}
%%%%%%%%%%%%%%%%%%%%%%%%%%%%%%%%%%%%%%%%%%%%%%%%%%%%%%%%%%%%%%%%%%%%%%%%%%%%%%%

See also Appendix G (Master Glossary).

\begin{table}[htbp]
\centering
\renewcommand{\arraystretch}{1.2}
\small
\begin{tabular}{@{}lll@{}}
\toprule
\textbf{Symbol} & \textbf{Name} & \textbf{Range} \\
\midrule
$\tau$ & Softmax temperature & $(0, \infty)$, typical: $0.5$--$2.0$ \\
$b_{ij}^{(r)}$ & Belief of agent $i$ about candidate $j$ at round $r$ & $[0,1]$, $\sum_j = 1$ \\
$\eta$ & Learning rate (belief update weight) & $[0,1]$, typical: $0.2$--$0.5$ \\
$\hat{v}_j^{(r)}$ & Observed vote share of candidate $j$ in round $r$ & $[0,1]$ \\
$\beta_{\text{thresh}}$ & Bandwagon activation threshold & $[0,1]$, typical: $0.25$--$0.40$ \\
$\alpha_1, \alpha_2$ & Bandwagon acceleration (linear, quadratic) & $\geq 0$ \\
$\theta_c$ & Commitment threshold & $[0,1]$, typical: $0.50$--$0.70$ \\
$\epsilon$ & Exploration probability for committed voters & $[0,1]$, typical: $0.05$--$0.15$ \\
$q^*$ & Supermajority threshold & $(0.5, 1]$, conclave: $2/3$ \\
$\delta$ & Candidate dropout threshold & $[0,1]$, typical: $0.03$--$0.10$ \\
$\phi_{\text{drop}}$ & Dropout decay factor & $(0,1)$, typical: $0.3$--$0.7$ \\
$N_{\text{MC}}$ & Monte Carlo replications & $\mathbb{N}$, typical: $5$--$100$ \\
$K$ & Number of historical validation cases & $\mathbb{N}$, conclave: $12$ \\
\bottomrule
\end{tabular}
\end{table}


%%%%%%%%%%%%%%%%%%%%%%%%%%%%%%%%%%%%%%%%%%%%%%%%%%%%%%%%%%%%%%%%%%%%%%%%%%%%%%%
\section{Critical Foundations and Objections}
\label{sec:bp-foundations}
%%%%%%%%%%%%%%%%%%%%%%%%%%%%%%%%%%%%%%%%%%%%%%%%%%%%%%%%%%%%%%%%%%%%%%%%%%%%%%%

\subsection{Objection 1: Why Not Full Bayesian Updating?}

\textbf{Objection:} DeGroot linear updating (VTM-2) is a simplification of true
Bayesian reasoning.

\textbf{Response:} Full Bayesian updating requires specifying a generative model
for votes, which introduces more free parameters than data can identify with $K=12$
cases. DeGroot updating converges to the same consensus
\citep{degroot1974reaching, chamley2004rational} and is empirically well-supported
for bounded-rational agents \citep{ambuhl2021belief}.

\subsection{Objection 2: Are Cascades Really Rational?}

\textbf{Objection:} Information cascades assume rational agents, but conclave voting
involves emotional, spiritual, and political factors.

\textbf{Response:} The cascade mechanism is \emph{robust} to bounded rationality.
\citet{bikhchandani1992theory} show that cascades form even with imperfect
rationality---what matters is that agents respond to public information (vote counts).
The QRE framework (VTM-1) already accommodates bounded rationality through $\tau > 0$.

\subsection{Objection 3: Overfitting with 8 Building Blocks?}

\textbf{Objection:} Eight mechanisms with multiple parameters may overfit 12 data points.

\textbf{Response:} Pattern-Oriented Modeling (VTM-8) addresses this by requiring
the model to simultaneously match \emph{three independent patterns} (winner, duration,
correlation), not just one. Moreover, the parameters are constrained by literature
priors (not freely estimated), and the Monte Carlo approach tests robustness.
\citet{grimm2005pattern} argue that multiple-pattern validation is a stronger test
than single-metric goodness-of-fit.


%%%%%%%%%%%%%%%%%%%%%%%%%%%%%%%%%%%%%%%%%%%%%%%%%%%%%%%%%%%%%%%%%%%%%%%%%%%%%%%
\section{References}
\label{sec:bp-references}
%%%%%%%%%%%%%%%%%%%%%%%%%%%%%%%%%%%%%%%%%%%%%%%%%%%%%%%%%%%%%%%%%%%%%%%%%%%%%%%

\nocite{bcm_master}

Key references for Appendix BP:

\begin{itemize}[nosep, leftmargin=*]
\item Axelrod, R. (1997). \textit{The Complexity of Cooperation}. Princeton University Press.
\item Banerjee, A.~V. (1992). A simple model of herd behavior. \textit{QJE}, 107(3), 797--817.
\item Bikhchandani, S., Hirshleifer, D., \& Welch, I. (1992). A theory of fads, fashion, custom, and cultural change as informational cascades. \textit{JPE}, 100(5), 992--1026.
\item Buchanan, J.~M., \& Tullock, G. (1962). \textit{The Calculus of Consent}. University of Michigan Press.
\item Caplin, A., \& Nalebuff, B. (1988). On 64\%-majority rule. \textit{Econometrica}, 56(4), 787--814.
\item Chamley, C. (2004). \textit{Rational Herds}. Cambridge University Press.
\item Crokidakis, N. (2025). When cardinals strategize: An agent-based model of influence and ideology for the papal conclave. \textit{arXiv:2505.07014}.
\item DeGroot, M.~H. (1974). Reaching a consensus. \textit{JASA}, 69(345), 118--121.
\item Epstein, J.~M. (2006). \textit{Generative Social Science}. Princeton University Press.
\item Galam, S., \& Jacobs, F. (2007). The role of inflexible minorities in the breaking of democratic opinion dynamics. \textit{Physica A}, 381, 366--376.
\item Grazzini, J., \& Richiardi, M. (2015). Estimation of ergodic ABMs by simulated minimum distance. \textit{JEDC}, 51, 148--165.
\item Grimm, V. et al.\ (2005). Pattern-oriented modeling of agent-based complex systems. \textit{Science}, 310(5750), 987--991.
\item Lux, T. (1995). Herd behaviour, bubbles and crashes. \textit{Economic Journal}, 105(431), 881--896.
\item Masuda, N. (2012). Evolution of cooperation driven by zealots. \textit{Scientific Reports}, 2, 646.
\item McFadden, D. (1974). Conditional logit analysis of qualitative choice behavior. In \textit{Frontiers in Econometrics}.
\item McKelvey, R.~D., \& Palfrey, T.~R. (1995). Quantal response equilibria for normal form games. \textit{G\&EB}, 10, 6--38.
\item Mobilia, M. (2003). Does a single zealot affect an infinite group of voters? \textit{PRL}, 91, 028701.
\item Robert, C.~P., \& Casella, G. (2004). \textit{Monte Carlo Statistical Methods}. Springer.
\item Train, K.~E. (2009). \textit{Discrete Choice Methods with Simulation}. Cambridge University Press.
\item Windrum, P., Fagiolo, G., \& Moneta, A. (2007). Empirical validation of agent-based models. \textit{JASSS}, 10(2), 8.
\end{itemize}
